\clearpage
\fancyhead[L]{\small VON LINDENAU. STERNWARTE SEEBERG.}
\hypertarget{native-503-lower}{}\label{native-503-lower}
\newcommand{\SeebergResidualMark}{\hbox{\vrule width0pt height3.272222bp depth0pt\special{pdf:bcontent}\special{pdf:code q 0 g 2.927778 1.377778 m 2.927778 1.205556 l 2.755556 1.033333 l 2.755556 0.172222 l 2.927778 0.000000 l 3.100000 0.172222 l 3.100000 1.377778 l h 0.172222 3.272222 m 0.000000 3.100000 l 0.000000 2.927778 l 0.172222 2.755556 l 1.205556 2.755556 l 1.377778 2.927778 l 2.066667 2.927778 l 2.238889 2.755556 l 2.755556 2.755556 l 2.927778 2.583333 l 2.927778 2.411111 l 3.100000 2.411111 l 3.100000 2.583333 l 3.444444 2.927778 l 6.027778 2.927778 l 6.200000 3.100000 l 6.027778 3.272222 l 3.961111 3.272222 l 3.788889 3.100000 l 3.616667 3.100000 l 3.444444 3.272222 l 2.583333 3.272222 l 2.411111 3.100000 l 0.516667 3.100000 l 0.344444 3.272222 l h f* Q}\special{pdf:econtent}\kern6.200000bp}}
\newcommand{\SeebergMoon}{\text{\mdseries\AstroSymbols ☽}}

\begin{center}
\small Göttingische gelehrte Anzeigen. No. 106. pp. 1049..1052. 2 July 1808.\\[3pt]
\rule{.72\linewidth}{.3pt}
\end{center}

We have recently received a reliable account of the present state of the famous \emph{Seeberg} Observatory from Mr \textsc{von Lindenau}, in a paper devoted to the subject. An excerpt from it will be all the more welcome to lovers of astronomy because for some time a number of unfounded reports about it have reached the public from dubious sources. The disruptions of war, together with some repairs that had become necessary on the roof of the observatory, made it necessary to take down the instruments in 1806. Various circumstances, which we need not list here, prevented the necessary arrangements for setting them up again until the beginning of this year. Mr \textsc{von Lindenau}, who had once before supervised the observatory in the absence of Mr \textsc{von Zach}, was entrusted with the task by His Serene Highness the Duke of Saxe-Gotha and took up residence at the observatory on 9 April. Bad weather delayed the start of the reinstallation until 20 April. Afterwards several clear nights favoured the work so much that by the end of April the transit instrument, the quadrant and the regulator clock were ready for observations. A fine series of numerous observations, made between 29 April and 9 May and communicated to us, shows how successfully Mr \textsc{von Lindenau} adjusted the instruments straight away. The limited space in these pages allows us to reproduce only a few of them.

\begin{center}
\emph{Observations of the Sun.}

\smallskip
{\small\setlength{\tabcolsep}{3pt}
\begin{tabular}{r|rrr|rrr|r}
1808 & \multicolumn{3}{c|}{\shortstack{Mean time\\at Seeberg}} & \multicolumn{3}{c|}{\shortstack{Right ascension\\of the Sun}} & \shortstack{Error of the\\tables} \\
\hline
April 30 & $23^{\mathrm{u}}$ & $57^{\mathrm{m}}$ & $4^{\mathrm{s}}16$ & $37^{\circ}$ & $29'$ & $4''2$ & $+3''8$ \\
May 3 & 23 & 56 & 42.06 & 40 & 20 & 57.4 & $+1.0$ \\
5 & 23 & 56 & 29.87 & 42 & 16 & 9.9 & $+2.5$ \\
6 & 23 & 56 & 24.89 & 43 & 14 & 3.3 & $-2.0$ \\
7 & 23 & 56 & 20.06 & 44 & 11 & 59.0 & $-0.5$ \\
9 & 23 & 56 & 11.73 & 46 & 8 & 15.5 & $\SeebergResidualMark\,2.6$ \\
\hline
\multicolumn{7}{r|}{Mean error of the tables} & $+1.8$ \\
\end{tabular}}
\end{center}

\clearpage
\SourceLeaf{504}
\begin{center}
\emph{Observations of the Moon.}

\smallskip
{\small\setlength{\tabcolsep}{3pt}
\begin{tabular}{r|rrr|rrr|r}
1808 & \multicolumn{3}{c|}{\shortstack{Mean time\\at Seeberg}} & \multicolumn{3}{c|}{\shortstack{Right ascension of the\\western $\SeebergMoon$ limb}} & \shortstack{Error of the\\tables.} \\
\hline
May 2 & $5^{\mathrm{u}}$ & $20^{\mathrm{m}}$ & $37^{\mathrm{s}}63$ & $120^{\circ}$ & $34'$ & $0''3$ & $+2''5$ \\
3 & 6 & 10 & 48.31 & 134 & 7 & 50.4 & $+0.6$ \\
4 & 7 & 0 & 59.58 & 147 & 41 & 52.8 & $+5.7$ \\
5 & 7 & 51 & 44.20 & 161 & 24 & 15.9 & $+6.2$ \\
6 & 8 & 43 & 44.35 & 175 & 25 & 30.6 & $+3.1$ \\
7 & 9 & 37 & 44.89 & 189 & 57 & 4.3 & $+2.3$ \\
8 & 10 & 13 & 19.15 & 205 & 7 & 5.8 & $+6.7$ \\
\hline
\multicolumn{7}{r|}{Mean error of the tables} & $+3''9$ \\
\end{tabular}}
\end{center}

The observations were compared with Mr \textsc{von Zach}’s new solar tables and \textsc{Bürg}’s lunar tables.

Mr \textsc{von Lindenau} also sent us 8 observed positions of \emph{Uranus}, 11 of \emph{Saturn}, 6 of \emph{Jupiter}, 9 of Venus and 8 of \emph{Mercury}. We shall give only the final results that Mr \textsc{von Lindenau} drew from them. The mean error of \textsc{de Lambre}’s Uranus tables at this epoch was $-0''8$ in longitude and $-16''$ in latitude. The mean error of the Saturn tables was $+9''1$ in longitude and $-4''3$ in latitude. Saturn was at opposition in 1808 on 9 May at $8^{\mathrm{u}}25^{\mathrm{m}}47^{\mathrm{s}}2$ mean time at Seeberg. The mean error of the Jupiter tables was $-1''6$ in longitude and $-1''9$ in latitude. Mr \textsc{von Lindenau} did not compare the observations of Venus and Mercury, because he found the tables of these planets no longer adequate to the present state of astronomy. For a year he has himself been collecting material for new tables of Venus, for which he has already gathered more than 200 good observations. With these and the equations of perturbation developed by \textsc{la Place}, he hopes soon to be able to produce something better than the existing tables.

Mr \textsc{von Lindenau} has made it the main object of his practical work during his stay at the Seeberg Observatory to determine refractions and the annual parallax of the fixed stars in right ascension. If there are grounds for believing that $\alpha$ Lyrae shows a parallax in declination of $4\ldots5''$, the parallax in right ascension must be $7''$, so that the difference between the positive and negative maxima rises to nearly one second of time. If such a star is compared with several small stars in its vicinity, one may therefore hope that a difference of this size will be detectable with the excellent 8-foot transit instrument by \textsc{Ramsden}. Admittedly, comparing it with small stars at the positive and negative maxima will present considerable \emph{practical} difficulties.

Mr \textsc{von Lindenau} has already begun observing refraction. Although it is still too early to draw a definite result, all his observations agree in giving, at a barometer reading of 28 inches and a temperature of $10^{\circ}$ Réaumur, a horizontal refraction of $37\ldots40'$. (We note on this occasion that Inspector \textsc{Bessel} too has concluded from a very careful examination of \textsc{Bradley}’s observations that the usual value for horizontal refraction must be increased considerably). — More generally, after a year Mr \textsc{von Lindenau} will give astronomers a detailed account of the results of his stay at the Seeberg Observatory.

The entire building that houses the observatory proper has now been fully restored. Work is under way on rebuilding the residence, which will hopefully also be completed in the course of this summer.
